% =====================================================================
% Paper B — Systems (MDPI), SI "Using Digital AI Systems as a Response
% to High Economic Turbulence and Uncertainty"
%
% Full first draft: abstract, narrative sections, and RQ1--RQ3 results written.
% Remaining [DRAFT] markers are back-matter only (CRediT, AI-use disclosure,
% Zenodo DOI) pending co-author confirmation. See ../plan/ and ../explore/.
% =====================================================================

\documentclass[systems,article,submit,moreauthors,pdftex]{Definitions/mdpi}

% --- Compatibility guards (no-ops under the current MDPI class) ---
\providecommand{\TitleCitation}[1]{}
\providecommand{\AuthorCitation}[1]{}
\providecommand{\orcidA}{}
\providecommand{\orcidB}{}
\providecommand{\orcidC}{}
\providecommand{\dataavailability}[1]{\vspace{6pt}\noindent\textbf{Data Availability Statement:} #1\par}
\providecommand{\institutionalreview}[1]{\vspace{6pt}\noindent\textbf{Institutional Review Board Statement:} #1\par}
\providecommand{\informedconsent}[1]{\vspace{6pt}\noindent\textbf{Informed Consent Statement:} #1\par}

% --- Extra packages used by the body (harmless if already provided by the class) ---
\usepackage{booktabs}
\usepackage{tabularx}
\usepackage{multirow}
\usepackage{amsmath}
\usepackage{amssymb}
\usepackage{xurl}

\graphicspath{{../figures/}}

% --- Article metadata (tentative) ---
\Title{Decision Fragility of a Fuzzy Multi-Objective Allocation System under Input Turbulence: Reproducibility and Operability as Guardrails}
\TitleCitation{Decision Fragility of a Fuzzy Multi-Objective Allocation System under Input Turbulence}

\newcommand{\orcidauthorA}{0000-0003-1534-6680}% Peyman Rabiei
\newcommand{\orcidauthorB}{0000-0003-0292-7435}% Daniel Arias-Aranda
\newcommand{\orcidauthorC}{0000-0002-1551-419X}% Vladimir Stantchev

\Author{Peyman Rabiei $^{1}$\orcidA{}, Daniel Arias-Aranda $^{1}$\orcidB{} and Vladimir Stantchev $^{2,}$*\orcidC{}}
\AuthorNames{Peyman Rabiei, Daniel Arias-Aranda and Vladimir Stantchev}
\AuthorCitation{Rabiei, P.; Arias-Aranda, D.; Stantchev, V.}

\address{%
$^{1}$ \quad Faculty of Economics and Business, University of Granada, Granada, Spain; prabiei@correo.ugr.es (P.R.); darias@ugr.es (D.A.-A.)\\
$^{2}$ \quad Institute of Information Systems, SRH University Heidelberg, Heidelberg, Germany; stantchev@computer.org}

\corres{Correspondence: stantchev@computer.org}

\abstract{%
Digital AI systems increasingly allocate scarce resources under turbulence and uncertainty,
yet whether such a system can be trusted when its inputs are degraded mid-shock is rarely
evaluated. We study a fuzzy multi-objective allocation system (Mamdani fuzzy inference with
NSGA-II/NRGA) as a digital decision system under input turbulence, on the acute
resource-allocation shock of post-disaster relief. Perturbing the raw input fields and scoring
every decision on the clean ground truth, we find the system is markedly \emph{input-fragile}:
modest degradation of the inputs its objectives are computed from re-routes most of the
allocation --- far more than a crisp baseline --- onto worse decisions, at two problem scales,
whereas turbulence in other inputs attenuates the effect. We argue this localised fragility
makes bit-for-bit reproducibility and auditability load-bearing trust guarantees rather than
niceties, and show the system is reproducible (within-environment REP $=1$, with an audit
trail) and operable (a full Pareto set in seconds using a few megabytes on commodity
hardware). The reference implementation is open source.}

\keyword{digital decision systems; multi-objective optimization; fuzzy inference systems; input uncertainty; reproducibility; auditability; decision stability; humanitarian logistics}

\begin{document}

% =====================================================================
\section{Introduction}\label{sec:intro}

High economic turbulence turns resource allocation into a decision made under two pressures at
once: scarcity, because a shock has pushed demand past supply, and degraded information,
because the data a decision rests on is itself disrupted while the shock unfolds. Digital AI
systems increasingly make these calls, yet the property that determines whether such a system
can be trusted in turbulence --- how its decisions behave when its inputs are degraded --- is
rarely evaluated. A system validated only for decision \emph{quality} on clean data can still
be unfit for turbulence if small input degradation makes it re-decide unpredictably.

We study this question on a fuzzy multi-objective allocation system: a Mamdani
fuzzy-inference front-end feeding NSGA-II/NRGA, which allocates scarce capacity across demand
points. Post-disaster relief allocation --- directing affected people to relief centres when
capacity is short and the situational data is disrupted --- is the acute instance we use: a
resource-allocation decision under an information-degrading shock. The system builds on the
authors' prior FIS-MOEA work for disaster operations~\cite{rabiei2021,rabiei2023}; a companion
study treats its resilience-theoretic formulation~\cite{paperA2026}, whereas here the object of
study is the \emph{system} --- its trustworthiness and deployability under input turbulence ---
rather than the optimisation theory.

We expected the fuzzy encoding to buy graceful degradation. It does not. The MOEA decision is
markedly input-fragile, and the fragility is localised to the inputs its objectives are
computed from; where turbulence hits other inputs, a crisp baseline is at least as fragile.
Because a single run's decision is this sensitive, we argue that bit-for-bit reproducibility
and auditability are the load-bearing trust guarantees for such a system, and we show the
decision is cheap enough to be deployable on commodity hardware. Concretely, the paper
contributes: (i) an input-turbulence fragility analysis of the fuzzy-MOEA decision against a
crisp baseline, at two problem scales; (ii) reproducibility and auditability as the guardrails
that fragility necessitates, with an operational reproducibility test; and (iii) an operability
envelope establishing deployability. Section~\ref{sec:background} sets the systems context;
Section~\ref{sec:system} describes the decision system; Section~\ref{sec:method} the protocols;
Section~\ref{sec:results} the results; and Sections~\ref{sec:discussion}--\ref{sec:conclusion}
the implications.

\section{Background}\label{sec:background}

We take a systems view of an AI decision aid: not only the algorithm that computes an
allocation, but the running artefact that ingests disrupted data, commits to a decision, and
must be trusted and operated under pressure. Three strands meet in that view.

\emph{Fuzzy inference under uncertainty.} Mamdani fuzzy inference systems encode expert
judgement as linguistic if--then rules over graded membership functions, and are attractive
precisely when inputs are imprecise or partially known --- as situational data is during a
shock. Combined with multi-objective evolutionary optimisation, they have been applied to
post-disaster operations, including vehicle routing and relief distribution~\cite{rabiei2021}
and volunteer-to-vacancy assignment~\cite{rabiei2023}.

\emph{Multi-objective evolutionary optimisation.} Algorithms such as NSGA-II and NRGA return a
Pareto front of non-dominated trade-offs rather than a single optimum, leaving the final
commitment to a decision-maker. This is a strength for transparency but, as we show, a source
of instability under input turbulence: the front, and the point committed from it, both move
when the objective landscape moves.

\emph{Decision stability, reproducibility, and operability.} Robustness of multi-objective
methods is usually studied against algorithmic stochasticity or against perturbation of
elicited preferences; the companion study~\cite{paperA2026} takes the latter route. Far less
attention is paid to stability of the committed decision under degradation of the \emph{input
data} itself, or to whether the running system is bit-for-bit reproducible and auditable ---
the properties that let a stakeholder trust a fragile decision at all. This paper addresses
that gap for a concrete digital allocation system.

\section{The Decision System}\label{sec:system}

The system allocates scarce relief capacity by directing a subset of affected people to
relief centres. It reads situational inputs --- per-person attributes (age, disability,
injury and living status, local infrastructure-damage level, and remaining resource time),
per-centre state (occupancy and resource-depletion rates), and per-route information (travel
duration, road condition, possible hazard) --- and converts them, through three Mamdani fuzzy
inference systems, into three objectives to be minimised: a fairness-in-prioritisation level,
a transport-infeasibility level, and a centre-allocation-imbalance level. The fuzzy front-end
is the design choice under test: it encodes expert judgement as linguistic rules that tolerate
imprecise inputs, rather than assuming precise ones.

A multi-objective evolutionary optimiser (NSGA-II, with NRGA as an alternative) searches the
space of allocations and returns a Pareto front of non-dominated trade-offs. Each candidate
is an integer-plus-real encoding that names which people are directed and to which centre;
every candidate the operators can produce is feasible by construction. The objective of any
allocation is computed at a single evaluation point against a pre-computed FIS cache, which
lets a decision made on one instance be scored against another's ground truth.

\textbf{Crisp baseline.} The foil is a deterministic greedy allocator that ranks people by a
crisp weighted sum of the same criteria and sends each to the crisp-cheapest centre, with no
fuzzy inference anywhere. It is scored through the identical evaluator, so any difference in
behaviour is attributable to the fuzzy front-end, not to a different objective.

\textbf{The tool.} The system is released as an open-source, MIT-licensed Python package
(\texttt{presidio-hardened-vol-assign}) with the experiment drivers and result manifests. It
is continuously tested and dependency-audited, and --- as Section~\ref{sec:rq2} establishes
--- reproducible bit-for-bit on stock hardware. We treat that production-grade,
audit-ready engineering as part of the contribution: a digital decision system intended for
use under turbulence must be one that others can run, check, and hold to account.

\section{Method}\label{sec:method}

\subsection{Instances and solver}\label{sec:method-setup}
We evaluate on two synthetic relief instances: \emph{small} (5 centres, 150 affected people,
50 directed) and \emph{large} (10 centres, 300 people, 100 directed). Both use the
three-objective model of Section~\ref{sec:system}, solved with NSGA-II (population 100,
150--200 generations); NRGA is included for the operability comparison. All runs are seeded;
determinism is a property we test rather than assume (Section~\ref{sec:rq2}).

\subsection{Input-turbulence protocol (RQ1)}\label{sec:method-turbulence}
Turbulence perturbs the raw input fields --- the situational data a planner collects
mid-shock --- not the elicitation weights. We perturb one field at a time in one of four
modes: additive Gaussian \emph{noise} (standard deviation a fraction of the field's plausible
range), \emph{bias} (a systematic shift), \emph{missingness} (a fraction of entries blanked
and imputed with the instance median or mode), and, for categorical fields, \emph{flip} (a
fraction relabelled). Each perturbation is applied at five levels ($0, 0.05, 0.1, 0.2, 0.4$;
level $0$ is an identity check) over twelve independent realisations, and every perturbed
value is clipped to its field's valid range so degradation cannot manufacture an
out-of-domain input.

The decision is \emph{made} on the degraded inputs but \emph{scored} on the clean ground
truth --- the realised consequence of deciding on bad data. Because the solver returns a
Pareto front rather than one allocation, we reduce each front to a single \emph{canonical
decision} (the solution minimising the equal-weight, min--max-normalised objective sum) and
compare it to the clean-input canonical decision. We report three stability metrics per
realisation: \emph{objective drift} (Euclidean distance between the two decisions' realised
objective vectors), \emph{allocation churn} (fraction of directed people re-routed), and
centre-load rank stability. For each realisation we regress a metric on the turbulence level
to obtain a degradation \emph{slope}, and test the fuzzy-vs-crisp slopes with a paired
Wilcoxon test. The crisp baseline (Section~\ref{sec:system}) decides by a deterministic
weighted-sum greedy rule on the same inputs and is scored through the identical evaluator, so
only the decision rule differs. We confirmed the fragility finding is not an artifact of the
canonical-decision rule by repeating it with a fixed-reference (nearest-ideal) rule
(Appendix / robustness).

\subsection{Reproducibility protocol (RQ2)}\label{sec:method-repro}
We hash each allocation front to a canonical, order-invariant SHA-256 signature over its
rounded objective vectors and person-to-centre assignments. Within-environment
reproducibility solves each (size, seed) twice and checks the signatures match; the signature
is recorded with an environment fingerprint (platform, Python, and library versions).
Cross-environment reproducibility compares the signature for matching (size, seed) across
target OS/Python builds, automated by the continuous-integration matrix.

\subsection{Operability measurements (RQ3)}\label{sec:method-operability}
We time one \texttt{solve()} call with a monotonic clock and capture its peak Python memory
with \texttt{tracemalloc}, excluding the one-off FIS pre-computation, over ten seeds per
(size, algorithm), and report median latency, peak memory, and throughput.

\section{Results}\label{sec:results}

\subsection{RQ1 --- Decision fragility under input turbulence}\label{sec:rq1}

We perturb one raw input field at a time (noise, missingness, or categorical
flip) across five turbulence levels, re-decide with both the fuzzy-MOEA system and
the crisp greedy baseline on each of twelve independent degradations, and score every
decision on the \emph{clean} ground truth. For each realisation we regress a stability
metric on the turbulence level to obtain a degradation \emph{slope}; Table~\ref{tab:rq1}
reports the median fuzzy and crisp slopes and a paired Wilcoxon test of the two.
We report the full eight-cell matrix on the small instance and a four-cell subset on the
large instance (10 centres, 300 people) to test scale-dependence.

\begin{table}[H]
\centering
\caption{RQ1 degradation slopes (stability metric vs.\ turbulence level), fuzzy-MOEA vs.\
crisp baseline, small instance, twelve realisations per cell. $p$ is a paired Wilcoxon
test of the per-realisation slopes.\label{tab:rq1}}
\small
\begin{tabular}{lrrrrrr}
\toprule
& \multicolumn{3}{c}{Allocation churn} & \multicolumn{3}{c}{Realised-objective drift}\\
\cmidrule(lr){2-4}\cmidrule(lr){5-7}
Field / mode & fuzzy & crisp & $p$ & fuzzy & crisp & $p$\\
\midrule
IDL / noise            & 1.28 & 0.49 & 0.0005 & 10.0 & 4.3 & 0.0005\\
IDL / missingness      & 1.27 & 0.16 & 0.0005 &  7.2 & 3.2 & 0.0005\\
Resource time / noise  & 1.28 & 0.42 & 0.0005 & 12.2 & 2.3 & 0.0005\\
Travel dur.\ / noise   & 1.30 & 0.62 & 0.0005 & 10.5 & 5.0 & 0.0005\\
Centre occ.\ / noise   & 1.25 & 0.73 & 0.0034 &  6.4 & 9.0 & 0.064\\
Centre occ.\ / missing & 2.06 & 0.00 & 0.0005 & 10.8 & 0.0 & 0.0005\\
Road cond.\ / flip     & 1.24 & 0.78 & 0.0005 &  4.8 & 7.8 & 0.009\\
Poss.\ hazard / flip   & 1.26 & 0.72 & 0.0005 &  5.4 & 4.3 & 0.151\\
\bottomrule
\end{tabular}
\end{table}

\textbf{Decision churn is far higher for the fuzzy-MOEA on the objective-coupled inputs.}
On the small instance the fuzzy churn slope exceeds the crisp baseline's in every one of
the eight cells ($p \le 0.0034$; Figure~\ref{fig:churn}). At scale
(Table~\ref{tab:rq1-large}) the gap stays sharp for the person-priority inputs the model's
objectives depend on --- infrastructure-damage-level noise and missingness give fuzzy churn
slopes $\approx 1.4$ against crisp $0.16$--$0.37$ ($p = 0.008$) --- but for centre-occupancy
noise the crisp baseline becomes just as churny (fuzzy $1.45$ vs.\ crisp $1.44$; $p = 0.95$).
The fragility is thus \emph{concentrated} where turbulence perturbs the inputs the optimiser
is chasing, rather than being a universal property of the solver.

\textbf{Realised-objective quality degrades for the same objective-coupled inputs.} The
churn translates into worse realised objectives chiefly for perturbations of the
person-priority and travel inputs (small: IDL, resource-time, travel-duration give a fuzzy
drift slope $2$--$3\times$ the crisp baseline, $p \le 0.0005$, Figure~\ref{fig:degradation-idl};
large IDL noise: fuzzy $7.5$ vs.\ crisp $2.4$, $p = 0.008$). For centre-occupancy noise the
crisp baseline is in fact the \emph{more} fragile on realised quality, markedly so at scale
(fuzzy $8.5$ vs.\ crisp $23.4$; $p = 0.016$) --- so both the churn and the quality advantages
of the crisp baseline are specific to inputs the MOEA objectives do not directly consume.

\begin{figure}[H]
\centering
\includegraphics[width=0.72\textwidth]{fig_churn_summary.pdf}
\caption{Allocation churn at turbulence level $0.2$, fuzzy-MOEA vs.\ crisp baseline, per
field/mode cell. The MOEA re-routes markedly more of the allocation under every turbulence
type.\label{fig:churn}}
\end{figure}

\begin{figure}[H]
\centering
\includegraphics[width=\textwidth]{fig_degradation_infrastructure_damage_level_noise.pdf}
\caption{Realised-objective drift (left) and allocation churn (right) vs.\ turbulence level
for infrastructure-damage-level noise; fuzzy-MOEA vs.\ crisp, mean $\pm 1$ s.d.\ The level-0
identity check is exact for both.\label{fig:degradation-idl}}
\end{figure}

\begin{table}[H]
\centering
\caption{RQ1 large-instance subset (10 centres, 300 people; 8 realisations). The fragility
persists for the objective-coupled person inputs (IDL) but attenuates for the
centre/transport inputs (centre-occupancy, road-condition).\label{tab:rq1-large}}
\small
\begin{tabular}{lrrrrrr}
\toprule
& \multicolumn{3}{c}{Allocation churn} & \multicolumn{3}{c}{Realised-objective drift}\\
\cmidrule(lr){2-4}\cmidrule(lr){5-7}
Field / mode & fuzzy & crisp & $p$ & fuzzy & crisp & $p$\\
\midrule
IDL / noise            & 1.44 & 0.37 & 0.008 & 7.5 &  2.4 & 0.008\\
IDL / missingness      & 1.42 & 0.16 & 0.008 & 5.4 &  4.8 & 0.55\\
Centre occ.\ / noise   & 1.45 & 1.44 & 0.95  & 8.5 & 23.4 & 0.016\\
Road cond.\ / flip     & 1.42 & 1.16 & 0.008 & 5.4 &  9.2 & 0.008\\
\bottomrule
\end{tabular}
\end{table}

The finding refines our starting hypothesis: fuzzy expert encoding does not buy graceful
degradation. The MOEA decision is markedly \emph{input-fragile}, but the fragility is
\emph{localised} to the inputs its objectives consume (person priority, travel), where both
the allocation and its realised quality degrade far more than a crisp baseline --- at both
problem scales. Where turbulence perturbs inputs the objectives do not directly use, the
crisp baseline's advantage narrows or reverses at scale. Either way, a system that
re-commits to a very different allocation under imperceptible input noise cannot be trusted
from a single run; that is precisely what makes the reproducibility and auditability
guarantees of Section~\ref{sec:rq2} load-bearing.

\subsection{RQ2 --- Reproducibility and auditability}\label{sec:rq2}

RQ1 shows a single run's allocation is input-fragile; the minimum trust requirement is then
that the run itself be reproducible and auditable. For each (size, seed) we solve twice
in-process and hash the resulting allocation front (a canonical, order-invariant SHA-256
over the objective vectors and person-to-centre assignments). All twenty runs reproduce
bit-for-bit (within-environment REP $= 1.0$), and each distinct seed yields a distinct
signature --- the hash is sensitive to the run, not a constant. Every run records an
environment fingerprint (platform, Python, and library versions) alongside its signature;
cross-environment reproducibility is then established by executing the same driver on each
target OS / Python and comparing the signature column, which the project's continuous-
integration matrix automates. Together with the run manifest, the dependency audit, and the
published security policy, the signature is an audit trail: a stakeholder can verify
\emph{which} run produced \emph{which} allocation after the fact. This reproducibility floor
does not remove the input fragility of RQ1 --- it makes the fragile decision accountable,
which is the precondition for trusting such a system under turbulence at all.

\subsection{RQ3 --- Operability envelope}\label{sec:rq3}

We profile the wall-clock latency and peak memory of one \texttt{solve()} call per
(size, algorithm) over ten seeds on commodity hardware (Table~\ref{tab:rq3}). The system
returns a full Pareto set in $4$--$9$ seconds using a few megabytes of working memory, at
both problem sizes; NRGA is consistently faster than NSGA-II. At $\sim$$400$--$830$
decisions per hour with a $\le 3$~MB footprint, the decision system is field-deployable on
a laptop without specialised infrastructure --- H-B3 is supported.

\begin{table}[H]
\centering
\caption{RQ3 operability envelope: median over ten seeds, 3-objective relief model, pop 100
/ gen 200.\label{tab:rq3}}
\small
\begin{tabular}{llrrr}
\toprule
Size & Algorithm & Latency (s) & Throughput (dec./h) & Peak memory (MB)\\
\midrule
Small & NSGA-II & 6.20 & 580 & 1.4\\
Small & NRGA    & 4.34 & 829 & 1.4\\
Large & NSGA-II & 8.99 & 400 & 2.6\\
Large & NRGA    & 6.78 & 531 & 2.6\\
\bottomrule
\end{tabular}
\end{table}

\section{Discussion}\label{sec:discussion}

Our starting expectation --- that a fuzzy front-end would let the system degrade gracefully
under input turbulence --- was wrong, and instructively so. The MOEA decision is
\emph{input-fragile}: modest degradation of the inputs its objectives are built from re-routes
most of the allocation, far more than a crisp baseline, and onto genuinely worse decisions.
The mechanism is direct: the optimiser chases objectives computed from the perturbed data, so
input noise reshapes the very landscape it searches. Where turbulence hits inputs the
objectives do not consume --- centre occupancy, road condition at scale --- the effect
attenuates and the crisp heuristic, which uses those inputs only weakly, is if anything the
more fragile of the two. Fragility is thus a property of the coupling between the perturbed
input and the objective, not of the solver in the abstract.

For AI-assisted decision-making under turbulence this carries two lessons. First, the inputs a
multi-objective system optimises over are exactly the inputs whose data quality must be
hardened before deployment; a system can look robust on held-out quality metrics yet re-decide
wildly when those specific inputs are noisy. Second, and more general: a decision system whose
single-run output is this sensitive to imperceptible input changes cannot be trusted from one
run. That is why we treat reproducibility and auditability not as software hygiene but as the
load-bearing trust guarantees (Section~\ref{sec:rq2}). They do not remove the fragility; they
make the fragile decision \emph{accountable} --- a stakeholder can establish which run produced
which allocation, and reproduce it --- which is the minimum precondition for using such a
system responsibly under pressure. The operability envelope (Section~\ref{sec:rq3}) shows that
this accountability comes at no practical cost: the system decides in seconds on a laptop.

\textbf{Limitations.} The study uses synthetic instances of a single application (post-disaster
relief, framed as an acute resource-allocation shock) and one solver family; the
economic-turbulence setting is engaged by analogy rather than with market data. The
input-fragility result is established at two problem sizes but on a four-cell large-instance
subset. Cross-environment reproducibility is defined and instrumented but is completed by the
continuous-integration matrix rather than reported here. These bound the claims to
allocation-style decision systems under input turbulence; whether the objective-coupling
mechanism generalises to other multi-objective decision systems is the natural next question.

\section{Conclusions}\label{sec:conclusion}

A fuzzy multi-objective allocation system, studied as a digital decision system under input
turbulence, is not robust in the way its expert-knowledge encoding suggests. Its committed
allocation is input-fragile --- degrading the inputs its objectives are built from re-routes
most of the allocation onto worse decisions, more so than a crisp baseline --- and the effect
holds at two problem scales, localised to those objective-coupled inputs. Rather than a defect
to hide, this reframes the trust question: because a single run is this sensitive, the
minimum requirement for deploying such a system under turbulence is that the run be
reproducible and auditable. We show it is (bit-for-bit within an environment, with an
audit trail), and that the decision is cheap enough --- seconds, a few megabytes --- to deploy
on commodity hardware. The practical takeaway is to harden the data quality of the specific
inputs a multi-objective system optimises over, and to make the system's decisions accountable
by construction. Whether the objective-coupling mechanism generalises to other multi-objective
decision systems under turbulence is the question we leave open.

% =====================================================================
\authorcontributions{\emph{[DRAFT — confirm with co-authors.]} Conceptualization, P.R., D.A.-A.\ and V.S.; software, P.R.\ and V.S.; formal analysis, V.S.; writing, P.R., D.A.-A.\ and V.S. All authors have read and agreed to the published version of the manuscript.}

\funding{This research received no external funding.}

\institutionalreview{Not applicable. This study used only synthetically generated problem instances and involved no human participants or animals.}

\informedconsent{Not applicable.}

\dataavailability{\emph{[DRAFT]} The open-source implementation, experiment drivers, and result manifests are available at \url{https://github.com/presidio-v/presidio-hardened-vol-assign}; a citable Zenodo release will be minted for this study.}

\acknowledgments{\emph{[DRAFT — AI-use disclosure to be confirmed with co-authors before submission.]}}

\conflictsofinterest{The authors declare no conflicts of interest.}

% =====================================================================
\reftitle{References}
\bibliography{lit}

\end{document}
